\subsection{Gas Composition Study\label{ch4:gasComposition}}
This section will investigate \gls{AO} properties for different gas composition
using a heat bath simulation. A heat bath simulation is an
excellent tool for examining the time it takes for a gas to reach equilibrium
after passing through a shock wave. The inputs for a heat bath simulation
typically include the flow characteristics observed after a shock wave. These
conditions are relatively straightforward to calculate in an experimental
setup, facilitating a more direct comparison between computational and
experimental results. \\

The solver used for this investigation is \gls{SU2}-\gls{NEMO} \cite{Palacios2013a}. \Cref{tab:gasCompositionStudyMassFraction}, shows the flow conditions for
cases studied during this analysis. These conditions were took from
\gls{hynecc} \cite{Gimelshein2022}. The \gls{hynecc} publication validates
\gls{SU2}-\gls{NEMO} by comparing the solver solution against other solvers
such as \todo{Look at hynecc paper and complete this}.
\begin{table}[H]
    \begin{center}
        \begin{tabular}{|c|c|c|c|c|}
        \hline
         Case & $T_{tr}\left[K\right]$ & $T_{rot}\left[K\right]$ &
            $T_{vib}\left[K\right]$ &  $P\left[kPa\right]$ \\ \hline\hline
            $1C$\cite{Gimelshein2022} & $20000$ & $20000$ & $300$ & $2760$ \\ \hline 
            $2C$\cite{Gimelshein2022} & $10000$ & $10000$ & $300$ & $1380$ \\ \hline 
            $3C$\cite{Gimelshein2022} & $15000$ & $15000$ & $300$ & $2067$ \\ \hline
        \end{tabular}
        \caption{Chemistry composition flow conditions}\label{tab:gasCompositionStudyMassFraction}
    \end{center}
\end{table}

    \Cref{tab:gasCompositionMassFractionConditions} shows the mass
    fraction composition for the cases in \Cref{tab:gasCompositionStudyMassFraction}. 

    \begin{table}[H]
        \centering
        \begin{tabular}{|c|c|c|c|c|c|}
            \hline
            Case  & $\ce{N}\Units{\,}$ & $\ce{O}\Units{\,}$ &
            $\ce{NO}\Units{\,}$ & $\ce{N_2}\Units{\,}$ &
            $\ce{O_2}\Units{\,}$ \\ \hline\hline
            $1C$ & $0.0$ & $0.0$ & $0.0$ & $1.0$ & $0.0$ \\ \hline
            $2C$ & $0.0$ & $0.0$ & $0.0$ & $0.0$ & $1.0$ \\ \hline
            $3C$ & $6.5E-7$ & $0.2283$ & $0.01026$ & $0.75$ & $0.713E\!-\!3$ \\ \hline
        \end{tabular}
        \caption{Mass fraction composition for cases in
        \Cref{tab:gasCompositionStudyMassFraction}}\label{tab:gasCompositionMassFractionConditions}
    \end{table}

    \subsubsection{Results}
    % Temperatures %
    \Cref{fig:temperaturesChemistryComposition} illustrates the translational and
    vibrational temperatures for cases described in \Cref{tab:gasCompositionStudyMassFraction}
    and \Cref{tab:gasCompositionMassFractionConditions}. As expected, the
    relaxation process takes longer for case $1C$, as it has the case with the
    highest initial translational temperature. In
    \Cref{fig:temperaturesChemistryComposition}, all cases reach thermal equilibrium at
    approximately $1\Units{\mu s}$. The diatomic nitrogen case illustrated in
    \Cref{fig:temperatures1C} displays a higher
    energy transfer between the translational and vibrational modes. 
    \begin{figure}[H]
        \centering
        \subfigure[Case
        $1C$]{\includegraphics[width=0.45\textwidth]{chapter4/gasComposition/1C_temperatures}}
        \subfigure[Case
        $2C$]{\includegraphics[width=0.45\textwidth]{chapter4/gasComposition/2C_temperatures}}
        \subfigure[Case
        $3C$]{\includegraphics[width=0.45\textwidth]{chapter4/gasComposition/3C_temperatures}}
        \caption{\label{fig:temperaturesChemistryComposition}Temperature results for cases in
            \Cref{tab:gasCompositionStudyMassFraction}}
    \end{figure}
% Temperatures %

    % Mass Fraction %%
    \Cref{fig:massFractionChemistryComposition} shows the mass fraction
    composition for the cases described in \Cref{tab:gasCompositionStudyMassFraction}
    and \Cref{tab:gasCompositionMassFractionConditions}.
    \Cref{fig:massFraction1C} shows the mass fraction for Case $1C$, which
    corresponds to a diatomic nitrogen gas. The diatomic nitrogen decomposes into nitrogen, and by
    the time $(\approx 10\Units{\mu s})$ the simulation reaches chemical equilibrium, approximately $20\%$ of the diatomic nitrogen has decomposed into nitrogen. 
    \Cref{fig:massFraction2C} shows the mass fraction for Case $2C$, which
    corresponds to diatomic oxygen gas.The diatomic oxygen decomposes into oxygen, and by
    the time $(\approx 10\Units{\mu s})$ the simulation reaches chemical
    equilibrium, approximately $20\%$ of the diatomic oxygen has decomposed into
    oxygen.
    \Cref{fig:massFraction3C} shows the mass fraction for Case $3C$,
    which corresponds to a $5$ species air mixture. Some of the diatomic nitrogen and oxygen
    decomposed to a nitric oxide, with a peak value of approximately $10\%$ at
    $5\Units{n s}$.
% Mass Fraction %%
    \begin{figure}[H]
        \centering
        \subfigure[Case
        $1C$]{\includegraphics[width=0.45\textwidth]{chapter4/gasComposition/1C_massFraction}}
        \subfigure[Case
        $2C$]{\includegraphics[width=0.45\textwidth]{chapter4/gasComposition/2C_massFraction}}
        \subfigure[Case
        $3C$]{\includegraphics[width=0.45\textwidth]{chapter4/gasComposition/3C_massFraction}}
        \caption{\label{fig:massFractionChemistryComposition}Mass fraction results for cases in
        \Cref{tab:gasCompositionStudyMassFraction}}
    \end{figure}
% Mass Fraction %%

    % Gladstone Species %
    \Cref{fig:speciesGladstoneDaleChemistryComposition} shows the species
    \gls{GD} constants for the cases described in
    \Cref{tab:gasCompositionStudyMassFraction} and
    \Cref{tab:gasCompositionMassFractionConditions}. These constants were
    calculated using \Cref{eq:speciesGladstoneDaleConst}. \Cref{fig:speciesGladstoneDale1C} shows the evolution of the diatomic
    nitrogen case. As expected, at the beginning of the simulation, the
    \gls{GD} constant for nitrogen is
    zero because the simulation contains only diatomic nitrogen.
    However, as the simulation evolves, some of the diatomic nitrogen gas
    dissociates
    into nitrogen causing the \gls{GD} constant for nitrogen to no longer be
    zero. \Cref{fig:speciesGladstoneDale2C} shows the evolution of the diatomic
    oxygen case. As expected, at the beginning of the simulation, the
    \gls{GD} constant for oxygen is
    zero because the simulation contains only diatomic oxygen.
    However, as the simulation evolves, some of the diatomic oxygen gas
    dissociates
    into oxygen causing the \gls{GD} constant for oxygen to no longer be
    zero. \Cref{fig:speciesGladstoneDale3C} illustrates the evolution of five
    gas species. Some of the diatomic nitrogen and oxygen recombine into nitric
    oxide, resulting in a peak in the \gls{GD} for nitric oxide at
    approximately $5\Units{ns}$. However, this peak is short-lived as the
    nitric oxide quickly dissociates back into nitrogen and oxygen. This figure
    highlights the direct impact of density changes on the \gls{GD}. As is well
    known, in a nonequilibrium flow, the high temperatures cause molecular
    dissociation. 
    \begin{figure}[H]
        \centering
        \subfigure[Case
        $1C$]{\includegraphics[width=0.45\textwidth]{chapter4/gasComposition/1C_speciesGladstoneDale}}
        \subfigure[Case
        $2C$]{\includegraphics[width=0.45\textwidth]{chapter4/gasComposition/2C_speciesGladstoneDale}}
        \subfigure[Case
        $3C$]{\includegraphics[width=0.45\textwidth]{chapter4/gasComposition/3C_speciesGladstoneDale}}
        \caption{\label{fig:speciesGladstoneDaleChemistryComposition}Species
        \gls{GD} results for cases in
        \Cref{tab:gasCompositionStudyMassFraction}}
    \end{figure}
    % Gladstone Species %

    % Gladstone Total %
    \Cref{fig:totalGladstoneDaleChemistryComposition} compares the total \gls{GD} between a nonequilibrium case and a frozen flow case for the conditions outlined in \Cref{tab:gasCompositionMassFractionConditions}. These results indicate that a higher percentage of diatomic nitrogen corresponds to an increased \gls{GD}. This outcome aligns with expectations, as diatomic nitrogen exhibits the highest equilibrium polarizability among the species considered. Additionally, it is noteworthy that nonequilibrium effects contribute to an increase in the \gls{GD}.
    \begin{figure}[H]
        \centering
        \subfigure[Case
        $1C$]{\includegraphics[width=0.45\textwidth]{chapter4/gasComposition/1C_totalGladstoneDale}}
        \subfigure[Case
        $2C$]{\includegraphics[width=0.45\textwidth]{chapter4/gasComposition/2C_totalGladstoneDale}}
        \subfigure[Case
        $3C$]{\includegraphics[width=0.45\textwidth]{chapter4/gasComposition/3C_totalGladstoneDale}}
        \caption{\label{fig:totalGladstoneDaleChemistryComposition}Total Gladstone-Dale results for cases in
        \Cref{tab:gasCompositionStudyMassFraction}}
    \end{figure}
    % Gladstone Total %

    % Refraction Index % 
    \Cref{fig:refractionIndexChemistryComposition} presents the index of refraction for the cases described in \Cref{tab:gasCompositionStudyMassFraction}. The figure compares the dilute formulation (\Cref{eq:diluteRefractiveIndex}) with the dense formulation (\Cref{eq:denseRefractiveIndex}) and \Cref{tab:gasCompositionMassFractionConditions}. It is observed that there is no significant difference between the dense and dilute formulations. As expected, the case shown in \Cref{fig:totalGladstoneDale1C}, which contains a higher proportion of diatomic nitrogen, exhibits a greater change in the refractive index.
    \begin{figure}[H]
        \centering
        \subfigure[Case
        $1C$]{\includegraphics[width=0.45\textwidth]{chapter4/gasComposition/1C_refractionIndex}}
        \subfigure[Case
        $2C$]{\includegraphics[width=0.45\textwidth]{chapter4/gasComposition/2C_refractionIndex}}
        \subfigure[Case
        $3C$]{\includegraphics[width=0.45\textwidth]{chapter4/gasComposition/3C_refractionIndex}}
    \caption{\label{fig:refractionIndexChemistryComposition}Index of Refraction results for cases in
    \Cref{tab:gasCompositionStudyMassFraction}}
    \end{figure}
% Refraction Index % 

    % Electric Susceptibility % 
    \Cref{fig:electricSusceptivilityResultsChemistryComposition} presents the
    electric susceptibility for the cases described in
    \Cref{tab:gasCompositionStudyMassFraction}. The results compares the
    dilute formulation with the dense formulation. It is observed that there is
    no significant difference between the dense and dilute formulations. As
    expected, the case shown in \Cref{fig:totalGladstoneDale1C}, which contains
    a higher proportion of diatomic nitrogen, exhibits a greater change in the
    electric susceptibility.
\begin{figure}[H]
\noindent
\hspace*{-0.08\textwidth}
\begin{minipage}{\textwidth}
    % First row: two subfigures
    \begin{minipage}[t]{0.39\textwidth}
        \centering
        \subfigure[Case $1C$]{
        \label{fig:electricSusceptibility1C}
        \includegraphics[height=\linewidth]{chapter4/gasComposition/1C_electricSusceptibility}
        }
    \end{minipage}
    \hspace{0.11\textwidth}  % small gap between the two figures
    \begin{minipage}[t]{0.39\textwidth}
        \centering
        \subfigure[Case $2C$]{
        \label{fig:electricSusceptibility2C}
        \includegraphics[height=\linewidth]{chapter4/gasComposition/2C_electricSusceptibility}
        }
    \end{minipage}
    \vspace{1em}
    % Second row: centered subfigure
    \begin{center}
    \begin{minipage}[t]{0.39\textwidth}
        \centering
        \subfigure[Case $3C$]{
        \label{fig:electricSusceptibility3C}
        \includegraphics[height=\linewidth]{chapter4/gasComposition/3C_electricSusceptibility}
        }
    \end{minipage}
    \end{center}
\end{minipage}
    \caption{\label{fig:electricSusceptivilityResultsChemistryComposition}Electric
    susceptibility results for cases in \Cref{tab:gasCompositionStudyMassFraction}}
\end{figure}
% Refraction Index % 
